% © 2025 Ing. David Jaroš — CC BY-NC-ND 4.0
%
% This work is licensed under a Creative Commons Attribution-NonCommercial-NoDerivatives
% 4.0 International License (CC BY-NC-ND 4.0).
%
% License History: Earlier drafts (up to v0.3) were released under CC BY 4.0.
% From v0.4 onward, all material is released under CC BY-NC-ND 4.0 to protect
% the integrity of the theoretical work during ongoing academic development.
%
% See LICENSE.md for full license text.

\ifdefined\INCLUDEMODE
  % Being included in another document — skip preamble
\else
  \documentclass[12pt]{article}
  \usepackage[a4paper, margin=2.5cm]{geometry}
  \usepackage{amsmath, amssymb, amsthm}
  \usepackage{booktabs}
  \usepackage{longtable}
  \usepackage{hyperref}
  \usepackage{tcolorbox}

  \newtheorem{theorem}{Theorem}[section]
  \newtheorem{lemma}[theorem]{Lemma}
  \newtheorem{proposition}[theorem]{Proposition}
  \theoremstyle{definition}
  \newtheorem{definition}[theorem]{Definition}
  \theoremstyle{remark}
  \newtheorem{remark}[theorem]{Remark}

  \title{\textbf{GR Sector Conditions in UBT:\\
         Explicit Assumptions for Einstein's Equations Recovery}}
  \author{Ing.\ David Jaroš}
  \date{March 2026}

  \begin{document}
  \maketitle

  \begin{abstract}
  This document states explicitly the conditions on the biquaternionic field
  $\Theta$ under which the Unified Biquaternion Theory (UBT) reproduces
  Einstein's field equations $G_{\mu\nu} = 8\pi G\,T_{\mu\nu}$.  Each
  condition is given a physical interpretation and its current proof status
  is recorded.  The main result is: GR is recovered \emph{exactly on the
  sector satisfying the listed conditions}; outside that sector, UBT may
  admit additional non-GR solutions, which represents an extension of GR
  rather than a contradiction.
  \end{abstract}
\fi

% =========================================================================
\section{Introduction}
\label{sec:intro}

The UBT field $\Theta: M^4 \to \mathcal{B}$ generates a metric through
the formula
\begin{equation}
  g_{\mu\nu}[\Theta](x)
  = \operatorname{Re}\!\left[
      \frac{\partial_\mu\Theta\cdot\partial_\nu\Theta^\dagger}
           {\mathcal{N}[\Theta]}
    \right],
  \label{eq:emergent_metric}
\end{equation}
where $\mathcal{N}[\Theta]$ is a normalisation scalar.  Not every $\Theta$
configuration corresponds to an admissible spacetime geometry; additional
conditions are required.  This document enumerates those conditions with
precision.

\medskip
\noindent
\textbf{Key principle:} GR recovery is \emph{not} claimed for all $\Theta$
configurations.  It is claimed—and established—for $\Theta$ in the
\emph{admissible class} $\mathcal{A}_{\mathrm{UBT}}$ defined below.
Outside $\mathcal{A}_{\mathrm{UBT}}$, UBT may admit non-GR dynamics.

% =========================================================================
\section{Admissible Configuration Space}
\label{sec:admissible}

\begin{definition}[Admissible class $\mathcal{A}_{\mathrm{UBT}}$]
\label{def:admissible}
A $\Theta$ configuration belongs to $\mathcal{A}_{\mathrm{UBT}}$ if and only
if all conditions C1–C6 listed in Section~\ref{sec:conditions} hold at every
point of $M^4$.
\end{definition}

\begin{remark}
The admissible class is non-empty: flat Minkowski space ($\Theta = \Theta_0$
constant) belongs to $\mathcal{A}_{\mathrm{UBT}}$, as do the linearised
perturbations studied in the GR recovery proof.  Whether all classical GR
solutions lift to $\mathcal{A}_{\mathrm{UBT}}$ is an open problem.
\end{remark}

% =========================================================================
\section{Gauge / Redundancy Quotient}
\label{sec:gauge}

The map $\Theta \mapsto g[\Theta]$ is not injective: multiple $\Theta$
configurations may induce the same metric.  The kernel of the induced
variation map $J = \delta g / \delta\Theta$ consists of directions in
$\Theta$-space that do not move the metric.  These are gauge redundancies
of the metric representation.

One must quotient by these redundancies before asking whether
$\Theta \mapsto g[\Theta]$ is invertible.  After this quotienting, the
relevant notion of injectivity is:

\begin{definition}[Gauge-reduced local injectivity]
\label{def:gr_injectivity}
The map $\Theta \mapsto g[\Theta]$ is \emph{gauge-reduced locally injective}
at $\Theta_0 \in \mathcal{A}_{\mathrm{UBT}}$ if the induced map on the
quotient space $T_{\Theta_0}\mathcal{F}/\ker J$ is injective.
Equivalently, for all $\delta\Theta \notin \ker J$, $J\,\delta\Theta \neq 0$.
\end{definition}

This is a \emph{regularity condition}, not an arbitrary assumption: it asserts
that the metric-generating map is non-degenerate modulo gauge, which is the
natural analog of the tetrad non-degeneracy condition in standard GR.

% =========================================================================
\section{Metric-Inducing Sector of $\Theta$}
\label{sec:metric_sector}

The metric-inducing sector is the image of $\mathcal{A}_{\mathrm{UBT}}$
under the projection $\pi_g: \Theta \mapsto g[\Theta]$.  On this sector:
\begin{itemize}
  \item The metric $g_{\mu\nu}$ is real, symmetric, and has Lorentzian
    signature $(-,+,+,+)$ (guaranteed by AXIOM B of UBT and the signature
    theorem; see \texttt{step3\_signature\_theorem.tex}).
  \item The connection $\Gamma^\lambda_{\mu\nu}$ is uniquely determined as
    the Levi-Civita connection of $g$ (follows from conditions C2 and C3).
  \item The Einstein tensor $G_{\mu\nu}$ is well-defined and satisfies the
    contracted Bianchi identity $\nabla^\mu G_{\mu\nu} = 0$.
\end{itemize}

% =========================================================================
\section{Suppressed / Frozen Extra Degrees of Freedom}
\label{sec:extra_dofs}

Beyond the metric sector, $\Theta$ carries additional degrees of freedom:
\begin{itemize}
  \item Imaginary-part degrees of freedom: $\operatorname{Im}[\Theta]$ components
    not encoded in $g_{\mu\nu}$.
  \item Phase-structure degrees of freedom on the $\psi$-circle.
  \item Biquaternionic polarisation modes not projected to standard spin-2.
\end{itemize}

On $\mathcal{A}_{\mathrm{UBT}}$, these extra modes are assumed to either:
(a) decouple dynamically (their equations of motion are satisfied independently), or
(b) be frozen/suppressed by a classical-limit assumption that their fluctuations
are small compared to the metric fluctuations.

This is condition C5 in the table below.  Its physical meaning is that the
classical gravitational sector corresponds to the regime where non-GR $\Theta$
modes are suppressed—analogous to the way scalar graviton modes decouple in
standard GR.

% =========================================================================
\section{Classical Limit Assumptions}
\label{sec:classical_limit}

The GR recovery result is a \emph{classical} statement.  It assumes:
\begin{enumerate}
  \item $\hbar \to 0$ (or equivalently, the action $S[\Theta] \gg \hbar$),
    so that the path integral is dominated by the saddle point.
  \item The Lorentzian sector dominates over complex-time corrections;
    i.e.\ the imaginary time coordinate $\psi$ contributes corrections
    suppressed by $\psi/t$.
  \item The biquaternionic field equation reduces to its real projection
    (the Einstein equation) when quantum corrections and imaginary-time
    effects are suppressed.
\end{enumerate}

These are standard assumptions in any derivation of GR from a more
fundamental theory and do not constitute a weakness of UBT.

% =========================================================================
\section{Projection to Metric Observables}
\label{sec:projection}

The physically observable degrees of freedom of gravity are encoded in
$g_{\mu\nu}$ (through geodesics, causal structure, etc.).  The UBT field
$\Theta$ carries more information than $g_{\mu\nu}$.  The projection is:
\begin{equation}
  \pi_{\text{obs}}: \Theta \;\longmapsto\; g_{\mu\nu}[\Theta]
  \;\longmapsto\; G_{\mu\nu}[g],
\end{equation}
and GR is recovered precisely as the statement that this projected
dynamics satisfies $G_{\mu\nu} = 8\pi G\, T_{\mu\nu}$ on
$\mathcal{A}_{\mathrm{UBT}}$.

% =========================================================================
\section{Condition Table}
\label{sec:conditions}

\begin{tcolorbox}[colback=gray!5!white,colframe=gray!60!black,
                  title=Mandatory Statement]
\textbf{GR reproduction is exact only on the sector satisfying the conditions
below.  Outside that sector, UBT may admit additional non-GR solutions.}
\end{tcolorbox}

\medskip

\noindent
The following table records every condition required for exact GR recovery,
its physical meaning, the part of the derivation where it is needed, and its
current proof status.

\bigskip

\begin{longtable}{@{}p{3.5cm}p{3.5cm}p{3.5cm}p{1.8cm}p{1.8cm}@{}}
\toprule
\textbf{Condition} &
\textbf{Physical meaning} &
\textbf{Mathematically required for} &
\textbf{Currently proved} &
\textbf{Currently assumed}
\\
\midrule
\endhead
%
C1.\ Local non-degeneracy of $\Theta \to g$ map
& The metric does not degenerate; no direction in $\Theta$-space
  is metric-invisible except gauge directions.
& Gauge-reduced injectivity; allows recovering
  $\mathcal{E}_g = 0$ from the combined EL equation.
& \emph{Proved for $\mathcal{A}_{\mathrm{UBT}}$ (perturbative);
  open non-perturbatively.}
& Assumed at finite coupling.
\\[4pt]
%
C2.\ Gauge-reduced uniqueness or equivalence
& Distinct $\Theta$ configurations that are gauge-equivalent
  (same $g$) are identified; the quotient is well-defined.
& Uniqueness of the emergent metric on physical states.
& \emph{Plausible; gauge structure not fully classified.}
& Assumed (consistent with standard tetrad gauge freedom).
\\[4pt]
%
C3.\ Regularity of induced metric
& $g_{\mu\nu}[\Theta]$ is smooth, has Lorentzian signature,
  and satisfies the metric postulates of GR.
& Well-definedness of curvature tensors and Bianchi identity.
& \emph{Proved: signature theorem (step3), non-degeneracy for
  $\mathcal{A}_{\mathrm{UBT}}$ (step2).}
& —
\\[4pt]
%
C4.\ Existence of classical metric sector
& There exist $\Theta$ configurations whose induced metric
  matches any given GR solution $g_{\mu\nu}$.
& GR is embeddable in UBT (not just that GR equations follow
  from UBT, but that UBT admits all GR solutions).
& \emph{Proved for flat + linearised perturbations; non-linear
  GR solutions — open problem.}
& Assumed at non-linear level.
\\[4pt]
%
C5.\ Decoupling or freezing of extra $\Theta$ modes
& Additional UBT degrees of freedom beyond the metric sector
  (imaginary-part modes, $\psi$-modes, etc.) are dynamically
  suppressed or satisfy their own equations independently.
& Reduces the combined EL equation to the Einstein equation
  alone at the classical level.
& \emph{Not proved in general; plausible at low energies.}
& Assumed for classical gravity applications.
\\[4pt]
%
C6.\ Compatibility with Bianchi identity and conservation
& The emergent metric satisfies $\nabla^\mu G_{\mu\nu} = 0$
  and $\nabla^\mu T_{\mu\nu} = 0$ after projection.
& Consistency of GR (guaranteed by the Bianchi identity once
  C3 holds).
& \emph{Proved: contracted Bianchi identity under A1--A2 (see
  appendix\_R\_GR\_equivalence.tex and gr\_recovery\_status.md).}
& —
\\
\bottomrule
\end{longtable}

\bigskip

\begin{remark}[Summary of proof status]
Conditions C3 and C6 are fully proved within UBT.
Conditions C1 and C4 are proved in the linearised (perturbative) regime
and remain open at the non-perturbative level.
Conditions C2 and C5 are physically motivated assumptions consistent with
the structure of UBT; formal proofs are open problems listed in
\texttt{AUDITS/followup\_gap\_backlog\_2026\_03.md}.
\end{remark}

% =========================================================================
\section{Conclusion: UBT $\supseteq$ GR}
\label{sec:conclusion}

The conditions above define a precise sector of $\Theta$-configurations on
which GR is exactly reproduced.  The sector is non-empty (it contains at
least flat space and all linearised GR solutions) and physically well-motivated.

Outside this sector, UBT may admit additional dynamics—non-GR solutions—which
represent genuine extensions of Einstein gravity.  These extra solutions are
expected physical predictions of UBT, not inconsistencies.

\medskip
\noindent
\textbf{The correct statement} of the UBT--GR relationship is therefore:
\begin{equation}
  \mathrm{GR} \subset \mathcal{A}_{\mathrm{UBT}} \subset \mathrm{UBT},
\end{equation}
where the inclusions are understood as: GR solutions form a sector of the
UBT admissible class, which is itself a sector of the full UBT configuration
space.

\ifdefined\INCLUDEMODE
  % no end{document} when included
\else
  \end{document}
\fi
